\scanpage{54}
trisection de $\mathcal L$, de ``somme'' nulle.

\uline{La} \deleted{\unsure{classification}} \inserted{donnée} \uline{d'un système} \deleted{\unsure{préelliptique|préelliptiques|privilégié}} \uline{\unsure{anharmonique|anharmoniques}} \inserted{\uline{\unsure{pointés|pointé|portés}}} \uline{sur $S$} est donc équivalente \uline{fonctoriellement} \uline{à} \uline{celle d'un} \uline{\unsure{faisceaux|faisceau}} \deleted{\unsure{en}} \uline{inversible} \deleted{\illegible{1}} $\underline L$ \uline{sur} $S$, \uline{et d'une trisection} $X'$ \inserted{\uline{non ramifiée}} \uline{de} $\mathcal L=\mathrm V(\underline L)$, \uline{de} \deleted{\unsure{projectée nulle}} ``\uline{somme}'' \uline{nulle} \deleted{\unsure{dans}} \deleted{$\underline L^{\otimes3}$}. Mais on peut aussi définir les \inserted{\unsure{autres}} fonctions symétriques élémentaires \unsure{pour cette trisection}

% slanted
\margin{un \unsure{revêt\textsuperscript{mt}} \deleted{\unsure{galoisien de}} non ramifié \deleted{\unsure{groupe $G$}} d'ordre 3, et une section de}

% archive sheet number: 53
